\begin{solapartat}
\punts{0,2} $\phi = \vec{B}\cdot\vec{A}$

\smallskip
\phantom{\punts{0,2}} $\phi = BA\cos\theta = BA\cos(\omega t)$

\smallskip
% DUBTE: a l'original «ω = 2π rad»; per a 10 voltes cada segon, ω = 20π rad/s (valor que s'usa a la línia següent).
\punts{0,2} $\omega = 2\pi\,rad$

\smallskip
\punts{0,2} $\varepsilon = -N\,\dfrac{d\phi}{dt} = NBA\omega\sin(\omega t)$

\smallskip
\punts{0,2} $\varepsilon = 200\times 1,25\times 10^{-2}\times 16\times 10^{-4}\times 20\pi\times\sin(20\pi\,t) = 0,251\,V\sin(20\pi\,t)$

\smallskip
% DUBTE: amb l'argument en radians, ε(1,28 s) = 0,251·sin(25,6π) = −0,239 V; el valor 0,247 V de la pauta
% correspon a calcular el sinus en graus (sin 80,4°).
\punts{0,2} $\varepsilon(t = 1,28\,s) = 0,247\,V$
\end{solapartat}

\begin{solapartat}
\punts{0,4} $\varepsilon_{\max} = 0,251\,V$

\smallskip
\punts{0,4} $\varepsilon_{ef} = \dfrac{\varepsilon_{\max}}{\sqrt{2}} = 0,177\,V$

% A l'original la gràfica és a la dreta, amb «0,2 p» en vermell al costat; s'ha esborrat del retall.
\medskip
\noindent\begin{minipage}{\linewidth}
\punts{0,2}
\figura[0.55]{sol_fig1.png}
\end{minipage}
\end{solapartat}
